\documentclass[12pt]{article}

%% Basic document formatting
\usepackage{amsmath, amsthm, amssymb, amsfonts}
\usepackage{mathtools}
\usepackage{mathrsfs}
\usepackage{xspace}
\usepackage{thmtools}
\usepackage{graphicx}
\usepackage{tikz-cd}
\usepackage{setspace}
\usepackage{fancyhdr}
\usepackage{titling}
\usepackage{hyperref}
\usepackage[left=0.4in,right=0.4in,top=1in,bottom=1in]{geometry}
\usepackage{float}
\usepackage{tabularx}
\usepackage{booktabs}
\usepackage[utf8]{inputenc}
\usepackage[english]{babel}
\usepackage{framed}
\usepackage[dvipsnames]{xcolor}
\usepackage{environ}
\usepackage{tcolorbox}
\tcbuselibrary{theorems,skins,breakable}

\usepackage{enumitem}
\setlist[enumerate]{leftmargin=*}
\setlist[enumerate,1]{labelindent=\parindent}
\setlist[enumerate,2]{labelindent=0pt}

% Blackboard Bold
\newcommand{\Z}{\mathbb{Z}}                 % integers
\newcommand{\Zpl}{\mathbb{Z}_{+}}           % positive integers
\newcommand{\Znn}{\mathbb{Z}_{\geq 0}}      % nonnegative integers. 
\newcommand{\Q}{\mathbb{Q}}                 % rationals
\newcommand{\Qpl}{\mathbb{Q}_{+}}           % positive Rationals
\newcommand{\R}{\mathbb{R}}                 % reals
\newcommand{\Rpl}{\mathbb{R}_{+}}           % positive Reals
\newcommand{\C}{\mathbb{C}}                 % complex numbers
\newcommand{\F}{\mathbb{F}}                 % field

% Words
\newcommand{\st}{\text{ such that }}
\newcommand{\wrt}{\text{ with respect to }}
\newcommand{\with}{\text{ with }}
\newcommand{\ie}{\text{, i.e. }}
\newcommand*{\ditto}{\text{---}\texttt{"}\text{---}}

% Operators
\newcommand{\abs}[1]{\left|#1\right|}                   % absolute value
\newcommand{\norm}[1]{\abs{\abs{#1}}}
\newcommand{\floor}[1]{\left\lfloor #1 \right\rfloor}   % floor
\newcommand{\ceil}[1]{\left\lceil #1 \right\rceil}      % ceiling
\newcommand{\powset}[1]{\mathscr{P}(#1)}

% Category Theory 
\renewcommand{\hom}{\text{Hom}}
\newcommand{\homspace}[3]{\hom_{#1}(#2, #3)}
\newcommand{\coproduct}[9]{
  \begin{tikzcd}[ampersand replacement=\&]
      \& \& #1 \arrow[dll, bend right, "#6"] \arrow[dl, "#5"] \\
   #4 \& #3 \arrow[l, "#9"] \\
      \& \& #2 \arrow[ull, bend left, "#8"] \arrow[ul, "#7"]
  \end{tikzcd}
}

\newcommand{\product}[9]{ 
  \begin{tikzcd}[ampersand replacement=\&]
      \& \& #3 \\
      #1 \arrow[r, "#7"] \arrow[urr, bend left, "#5"] \arrow[drr, bend right, "#6"] \& #2 \arrow[ur, "#8"] \arrow[dr, "#9"] \\ 
      \& \& #4
  \end{tikzcd}
}


% Algebra 
\newcommand{\Syl}[2]{\text{Syl}_{#1}{(#2)}}           % set of Sylow-p subgroups 
\newcommand{\<}{\langle}                % \<x\>, subgroup generated by x 
\renewcommand{\>}{\rangle}
\renewcommand{\index}[2]{[#1 : #2]}
\newcommand{\id}{\text{id}}             % identity element
\newcommand{\lhdneq}{%
  \mathrel{\ooalign{$\lneq$\cr\raise.22ex\hbox{$\lhd$}\cr}}}
\newcommand{\order}[1]{\text{o}(#1)}    % order of an element
\newcommand{\spec}{\operatorname{Spec}}
\newcommand{\rad}{\operatorname{rad}}   % radical of an ideal 
\newcommand{\sym}[1]{\mathfrak{S}_{#1}}
\newcommand{\aut}[1]{\text{Aut}(#1)} 
\newcommand{\gal}[2]{\text{Gal}(#1 / #2)} 
\newcommand{\inn}[1]{\text{Inn}(#1)} 
\let\oldcong\cong
\let\oldequiv\equiv
\renewcommand{\cong}{\oldequiv}
\renewcommand{\equiv}{\oldcong}

\newcommand{\triangleleftneq}{\mathrel{\ooalign{%
  $\trianglelefteq$\cr
  \hidewidth\raise0.06ex\hbox{$\not\mkern4mu$}\hidewidth\cr
}}}

\makeatletter % cycle
\newcommand{\cyc}[1]{(\mathbf{\cyc@process#1\relax})}
\def\cyc@process#1#2\relax{%
	#1%
	\ifx\relax#2\relax
	\else
		\,\cyc@process#2\relax
	\fi
}
\makeatother

\newcommand{\GL}[2]{\text{GL}_{#1}(#2)}                             % general linear group
\newcommand{\SL}[2]{\text{SL}_{#1}(#2)}                             % special linear group
\newcommand{\gl}[2]{\mathfrak{gl}_{#1}(#2)}
\renewcommand{\sl}[2]{\mathfrak{sl}_{#1}(#2)}
\newcommand{\so}[2]{\mathfrak{so}_{#1}(#2)}
\newcommand{\su}[1]{\mathfrak{su}(#1)}

% Linear Algebra
\newcommand{\bmat}[1]{\begin{bmatrix}#1\end{bmatrix}}   % bracketed matrix
\newcommand{\rank}{\operatorname{rank}}                 % rank
\newcommand{\nullity}{\operatorname{nullity}}           % nullity
\newcommand{\tr}{\operatorname{tr}}

% Combinatorics
\newcommand{\qnum}[1]{\left[#1\right]_q}                % q-analog
\newcommand{\qfact}[1]{\left[#1\right]!_q}              % q-analog factorial 
\newcommand{\qbinom}[2]{\binom{#1}{#2}_q}               % Gaussian binomial coefficient

% Topology/Analysis
\newcommand{\ball}[2]{\text{B}_{#1}(#2)}  % B_r(x): open r-balls around x
\newcommand{\diam}{\text{diam}}           % diamter of a set in metric space 
\newcommand{\lowint}[2]{\underline{\int_{#1}^{#2}}}
\newcommand{\upint}[2]{\overline{\int}_{#1}^{#2}}

% Misc Notation
\newcommand{\oneton}[1]{\{1, \ldots, #1\}}
\newcommand{\defeq}{\vcentcolon=}   % :=
\newcommand{\eqdef}{=\vcentcolon}   % =: 
\renewcommand{\bf}[1]{\textbf{#1}}
\newcommand{\im}{\operatorname{im}}
\newcommand{\fnrest}[2]{\left.#1\right|_{#2}}
\newcommand{\isom}{\overset{\sim}{\rightarrow}} 


\renewcommand{\thesection}{\arabic{section}.}
\renewcommand{\thesubsection}{(\alph{subsection})}

\newtheorem{claim}{Claim}
\newtheorem*{lemma}{Lemma}

%% Headers & title setup
\newcommand{\course}{Math100C}
\newcommand{\myname}{Jay Ser}
\setlength{\headheight}{14.5pt}
\pagestyle{fancy}
\fancyhf{}
\renewcommand{\headrulewidth}{0.4pt}
\lhead{\course}
\rhead{\myname}
\cfoot{\thepage}
\setlength{\droptitle}{-4em} 
\title{\course\ - WK2} 
\author{\myname}
\date{2026.04.10}

\begin{document}
\maketitle
\thispagestyle{fancy}

%------------------------------------------------------------------------------%

%TODO: do I need to verify more things about the matrix to show that the representation is properly defined? 
\section{Standard Two-Dimensional Representation of $D_n$}

Write the standard two-dimensional representation of the dihedral group $D_n$
explicitly in terms of matrices with respect to some basis.

\noindent\dotfill

Let's just use the standard basis for $\R^2$. 
Denoting $c = \cos{\frac{2\pi}{n}}$ and $s = \sin{\frac{2\pi}{n}}$, recall that 
\begin{align*}
  A_x = \bmat{c & -s \\ s & c} & \text{ rotates any vector } v \in \R^2 \text{ by } \frac{2\pi}{n} \text{ around the origin} \\ 
  A_y = \bmat{1 & 0 \\ 0 & -1} & \text{ reflects any vector } v \in R^2 \text{ across the } e_1 \text{-axis}
\end{align*}
Letting $x = \text{rotation by } 2\pi / n$ and $y = \text{ reflection across } e_1$, the standard two-dimensional representation of $D_n$ is defined by 
\begin{align*}
  x & \mapsto A_x \\ 
  y & \mapsto A_y
\end{align*}
since $D_n$ is presented by $(x, y \mid x^n, y^2, xyxy)$. 

%------------------------------------------------------------------------------%

\section{Irreducible Representations of $\mathfrak{S}_3$ Are Isomorphic}
\newcommand{\grp}{\mathfrak{S}_3}

Let $R_1, R_2 \colon \mathfrak{S}_3 \to \mathrm{GL}_2(\mathbb{C})$ be two irreducible
representations.  Prove directly (without using any theorems about characters)
that $R_1$ and $R_2$ are isomorphic.

\emph{Hint:} let $x = (123)$ and $y = (12)$.  Pick a nonzero
$u \in \mathbb{C}^2$ and look at the $G$-orbits of
\[
  u + R_1(g)(u) + R_1(g^2)(u)
  \qquad\text{and}\qquad
  u + R_1(h)(u).
\]
\noindent\dotfill

Since I'm going to be changing basis a lot, I will let the two given representations be independent of basis: 
take vector space representations $R, S: \mathfrak{S}_3 \rightarrow \GL(\C^2)$. 
More so, in an abuse of notation, I will continue to use the symbols $R$ and $S$ after I fix some basis for $\C^2$, thus viewing $R$ and $S$ as matrix representations. 

I first consider the representation $R$. 
Let $x = \cyc{123}, \, y = \cyc{12}$. 
For notational clarity, I denote $R_x$ as the invertible linear transformation $R(x)$. 
$\grp$ is a finite group, so $R_{x}$ is diagonalizable. 
Let $\lambda$ be an eigenvector of $R_x$ with eigenvector $v_1$. 
Particularly, one can choose $\lambda \neq 1$: 
that is, at least one of the eigenvalues of $R_x$ is not 1. 
If $\lambda = 1$ is a eigenvalue of $R_x$ with multiplicity 2, then with respect to the eigenbasis, $R_x$ is a scalar matrix (namely, the identity). 
Scalar matrices commute in $\GL_2(\C)$; 
hence $R_x$ commutes with the other generator for $R(\grp)$, which is $R_y$ (since $\grp = \<x, y\>$),
i.e., $R(\grp)$ is an abelian subgroup of $\GL_{2}(\C)$. 
Factoring the map $R: \grp \rightarrow \GL_{2}(\C)$ into 
$$ \grp \rightarrow R(\grp) \overset{\iota}{\hookrightarrow} \GL_{2}(\C), $$
one sees that $R$ is the direct sum of irreducible representations because $\iota$ is a two-dimensional representation of the abelian group $R(\grp)$ (recall that all irreducible representations of abelian groups are degree 1). 
This contradicts the assumption that $R$ is an irredubible two-dimensional representation of $\grp$. 

To conclude, we can assume $\lambda \neq 1$ is an eigenvalue of $R_x$ with eigenvector $v_1$. 
Next, use the relation $yx = x^2 y$ in $\grp$ to calculate
\begin{align*}
  R_x R_y v_1 
  & = R_{xy} v_1 = R_{yx^2} v_1 = R_{y} R_{x} R_{x} v_1 \\ 
  & = R_{y} \lambda^2 v_1 = \lambda^2 R_{y} v_1.
\end{align*}
Namely, this shows $v_2 \defeq R_y v_1$ is an eigenvector of $R_x$ with eigenvalue $\lambda^2$. 
Recall that because $x$ is an order 3 element of $\grp$, $\lambda$ is a third root of unity. 
Since $\lambda \neq 1$, that is, $\lambda$ is a nontrivial third root of unity, $\lambda^2$ and $\lambda$ are distinct eigenvalues of $R_x$, so $v_1$ and $v_2$ are linearly independent vectors in $\C^2$. 
$(v_1, v_2)$ gives a basis for $\C^2$, with respect to whom 
$$R_x = \bmat{\lambda & \\ & \lambda^2}, \: R_y = \bmat{ & 1 \\ 1}.$$
The matrix $R_y$ comes from the fact that $v_2 = R_y v_1$ and, because $y$ is an order 2 element in $\grp$, $R_y v_2 = v_1$. 

Using the exact same line of reasoning for the representation $S$, find $w_1, w_2 \in \C^2$ such that $(w_1, w_2)$ is a basis for $\C^2$ and 
$$S_x = \bmat{\lambda \\ & \lambda^2}, \: S_y = \bmat{ & 1 \\ 1}$$
with respect to this basis. 
Notice that the two distinct nontrivial third roots of unity must be eigenvalues for $S_x$ as well: 
as was reasoned above, at least one of the eigenvalues for $S_x$ is a nontrivial third root of unity, and that eigenvalue gives the other nontrivial third root of unity as an eigenvalue for $S_x$. 

Now, define the invertible linear map 
$$T: \C^2 \rightarrow \C^2; \: v_1 \mapsto w_1, \, v_2 \mapsto w_2.$$
To verify that this is an isomorphism of the representations $R$ and $S$, the following commuting diagram must commute for all $g \in \grp$: 
\begin{center}
  \begin{tikzcd}
    \C^2 \arrow[d, "T"] \arrow[r, "R_g"] & \C^2 \arrow[d, "T"]\\ 
    \C^2  \arrow[r, "S_g"] & \C^2
  \end{tikzcd}
\end{center}
Particularly, since $x, y$ generate $\grp$ and $(v_1, v_2)$ is a basis for $\C^2$, one only needs to check the diagram commute for four cases: 
independently pick group elements $g \in \{x, y\}$ and vectors $v \in \{v_1, v_2\}$. 
And indeed, because the matrices $R_g = S_g$ for all $g \in \grp$, where $R_g$ is written with respect to $(v_1, v_2)$ and $S_g$ is written with respect to $(w_1, w_2)$, the diagram does commute. 
For example, for $g = x$ and $v = v_1$, 
\begin{center}
  \begin{tikzcd}
    v_1 \arrow[d, "T", maps to] \arrow[r, "R_g", maps to] & \lambda v_1 \arrow[d, "T", maps to]\\ 
    w_1  \arrow[r, "S_g", maps to] & \lambda w_1
  \end{tikzcd}
\end{center}

%------------------------------------------------------------------------------%

\section{Character Tables}

Compute (with justification) the full character tables for each of the
following groups.

\subsection{The Klein Four-Group}
Denote $V_4 = \{1, x, y, xy\}$ as the Klein four-group. 
Recall $V_4 \equiv C_2 \times C_2$; 
namely, $V_4$ is abelian, so each element forms its own conjugacy class, and it follows that there are four irreducible nonisomorphic representations of $V_4$. 
If $d_1, \ldots, d_4$ are the dimensions of the irreducible representations $\rho_1, \ldots, \rho_4$, then 
$$d_{1}^{2} + \ldots + d_{4}^{2} = 4$$ 
must be satisfied. 
Clearly, each $d_i = 1$. 
Each nonidentity element has order 2 in $V_4$, so each group operator $\rho_{z}$ for $z \in V_4$ must have trace = -1 or 1. 
For each $i$, notice that because $\rho_{i}$ is a one-dimensional representation, $\chi_{i}(x) \chi_{i}(y) = \chi_{i}(xy)$; 
the trace of $\rho_{i}(z)$ is equal to just the singular entry of the $1 \times 1$ matrix $\rho_{i}(z)$, and $\rho$ itself is a homomorphism.
Hence for each $i$, one can just make two binary choices of picking $\chi_{i}(x) \in \{1, -1\}$ and $\chi_{i}(y) \in \{-1, 1\}$ to fully determine $\chi_{i}$. 
This yields exactly the four possible characters of an irreducible representation of $V_4$. 
\begin{center}
  \begin{tabular}{l | cccc}
    $V_4$ & $[1]$ & $[x]$ & $[y]$ & $[xy]$ \\
    \midrule
    $\chi_1$ & 1 & 1  & 1 & 1 \\
    $\chi_2$ & 1 & 1 & -1 & -1 \\
    $\chi_3$ & 1 & -1 & 1 & -1 \\ 
    $\chi_4$ & 1 & -1 & -1 & 1
  \end{tabular}
\end{center}

\subsection{The Quaternion Group of Order $8$}
Denote the quaternion group $Q_8 = \{1, -1, i, -i, j, -j, k, -k\}$. 
Considering the presentation 
$$Q_8 = (-1, i, j, k \mid i^2 = 1, -i^2 = -j^2 = -k^2 = -ijk),$$
it is clear that the center of $Q_8$ is $\{1, -1\}$. 
Given the symmetric properties satisfied by $i, j$ and $k$, one might guess that the class equation and conjugacy classes for $Q_8$ are
$$8 = 1 + 1 + 2 + 2 + 2; \hspace{1.5em} \{1\}, \{-1\}, \{i, -i\}, \{j, -j\}, \{k, -k\}$$ 
Indeed, this is true. 
To give an example for the conjugacy class of $i$, $jij^{-1} = ji(-j) = -jk = -i$ and $kik^{-1} = -i$ as well. 

So there are five isomorphism classes for the irreducible representations of $Q_8$. 
Letting $d_1, \ldots, d_5$ denote the dimensions of nonisomprhic irreducible representations $\rho_1, \ldots, \rho_5$, 
\begin{equation} \label{eq:3_a_dimrep}
d_{1}^2 + \ldots + d_{5}^{2} = 8.
\end{equation}
Up to reordering, only 
$$(d_1, d_2, d_3, d_4, d_5) = (1, 1, 1, 1, 2)$$ 
satisfies \eqref{eq:3_a_dimrep}. 
Consider the table
\begin{center}
  \begin{tabular}{l | ccccc}
    $Q_8$ & $[1]$ & [-1] & $[i]$ & $[j]$ & $[k]$ \\
    \midrule
    $\chi_1$ & 1 & 1  & 1 & 1 & 1 \\
    $\chi_2$ & 1 & a & b & c & d \\
    $\chi_3$ & 1 &  &  & & \\ 
    $\chi_4$ & 1 &  &  &  & \\
    $\chi_5$ & 2 &  &  &  &
  \end{tabular}
\end{center}
Because $0 = \<\chi_1, \chi_2\> = \frac{1}{\#Q_8}\sum_{g \in Q_8} \overline{\chi_{1}(g)}\chi_{2}(g)$,
$$0 = 1 + a + 2b + 2c + 2d.$$
Since $\chi_2$ is a 1-dimensional representation, given $g \in Q_8$ of order $n$, $\chi_{2}(g)$ is an $n$th root of unity. 
Namely, $a = \pm 1$ and $b, c, d$ are fourth roots of unity. 
Finally, due to the relation $i^2 = j^2 = k^2 = ijk = -1$ in $Q_8$, $a = bcd$ and $a = b^2 = c^2 = d^2$. 
I reorganize these conditions below: 
\begin{enumerate}
  \item $0 = 1 + a + 2b + 2c + 2d.$
  \item The roots of unity condition.
  \item $a = bcd$.
  \item $a = b^2 = c^2 = d^2$. 
\end{enumerate}
If $a = -1$, then condition 1 says $b + c + d = 0$. 
But this is impossible given that $b, c, d$ are fourth roots of unity. 
So $a = 1$ is forced. 
By condition 4, $b, c, d \in \{\pm 1\}$. 
Since $\chi_2$ is not the trivial representation, at least one of $b, c, d$ is $-1$. 
Condition 3 then supplies further restriction; 
exactly two of $b, c, d$ must be $-1$. 
It follows that $(b, c, d) \in \{(1, -1, -1), (-1, 1, -1), (-1, -1, 1)\}$. 
In fact, this exact conclusion applies to $\chi_3$ and $\chi_4$ as well. 
Thus the following character table is forced for the 1-dimensional representations: 
\begin{center}
  \begin{tabular}{l | ccccc}
    $Q_8$ & $[1]$ & [-1] & $[i]$ & $[j]$ & $[k]$ \\
    \midrule
    $\chi_1$ & 1 & 1  & 1 & 1 & 1 \\
    $\chi_2$ & 1 & 1 & 1 & 1 & -1 \\
    $\chi_3$ & 1 & 1 & 1 & -1 & 1 \\
    $\chi_4$ & 1 & 1 & -1 & 1 & 1 \\
    $\chi_5$ & 2 &  &  &  &
  \end{tabular}
\end{center}
Since the matrix given by the character table is unitary after scaling, its columns must be orthogonal. 
This gives the final row of the table: 
\begin{center}
  \begin{tabular}{l | ccccc}
    $Q_8$ & $[1]$ & [-1] & $[i]$ & $[j]$ & $[k]$ \\
    \midrule
    $\chi_1$ & 1 & 1  & 1 & 1 & 1 \\
    $\chi_2$ & 1 & 1 & 1 & 1 & -1 \\
    $\chi_3$ & 1 & 1 & 1 & -1 & 1 \\
    $\chi_4$ & 1 & 1 & -1 & 1 & 1 \\
    $\chi_5$ & 2 & -2 & -1 & -1 & -1
  \end{tabular}
\end{center}


%------------------------------------------------------------------------------%

\section{Space of Linear Transformations as a Representation}
Let $R_1: G \rightarrow \GL(V_1)$ and $R_2: G \rightarrow \GL(V_2)$ be two representations of a group $G$. 
Let $V$ be the space of linear transformations $T \colon V_1 \to V_2$.

\subsection{Vector Space Structure on $V$}
\newcommand{\calH}{\mathcal{H}}

There is a ``natural'' way to view $V$ as a vector space over $\mathbb{C}$.
What are the addition and scalar multiplication?

\noindent\dotfill

I refuse to denote the space of linear transformations $V_1 \rightarrow V_2$ as $V$. 
Instead, I call it $\calH$. 
Given $T_1, T_2 \in \calH$, define the element $T_1 + T_2$ as the linear transformation 
$$(T_1 + T_2)(v) = T_1(v) + T_2(v)$$ 
for all $v \in V_1$. 
It is trivial to check that such a mapping is indeed a linear transformation. 
Next, for any $c \in \C$ and $T \in \calH$, let $cT$ be the linear transformation 
$$(cT)(v) = cT(v)$$
for all $v \in V_1$. 
It follows that $\calH$ can be viewed as a vector space over $\C$. 

\subsection{A Representation on $V$}

Show that with the vector space structure from (a), we get a representation
$R \colon G \to \mathrm{GL}(V)$ by the formula
\[
  R(g)(T)(v_1) \;=\; R_2(g)\!\left(T\!\left(R_1(g^{-1})(v_1)\right)\right).
\]
\noindent\dotfill 

$R$ is a map $G \rightarrow \GL(\calH)$, which is confusing to grasp. 
One needs to stare at the following derivation for a while before it makes sense. 
My advice is to just type-check each function-composition and expression. 
Take any $v \in V_1$ and any $T: V_1 \rightarrow T_2$. 
\begin{align*}
  \{[R(g_1)R(g_2)](T)\} (v) 
  & = \{[R(g_1)](\overbrace{[R(g_2)](T))}^{V_1 \rightarrow V_2} \} (v) \\ 
  & = R_{2} (g_1) \left\{ R(g_2) [T R_{1}(g_{1}^{-1})] (v) \right\}  \\ 
  & = R_{2} (g_1) \left\{ R_{2} (g_2) [T R_1(g_{1}^{-1})] R_{1}(g_{2}^{-1}) (v) \right \} \\ 
  & = R_2(g_1 g_2) T R_1(g_2^{-1} g_{1}^{-1}) (v) \\ 
  & = \left\{[R(g_1 g_2)] (T) \right\} (v)
\end{align*}
where the last equality follows because $R_1$ and $R_2$ are group representations. 
The manipulation above shows that $R: G \rightarrow \GL(\calH)$ is a homomorphism as well, which means $R$ is a group representation. 

%------------------------------------------------------------------------------%

\section{Finite Subgroups of $\mathrm{GL}_n(\mathbb{R})$ Are Conjugate into $O_n$}

By imitating the proof of Maschke's theorem, prove that every finite subgroup
$G$ of $\mathrm{GL}_n(\mathbb{R})$ is conjugate to a subgroup of the orthogonal
group $O_n$.

\emph{Hint:} construct a $G$-invariant positive definite inner product on $V$.

\noindent\dotfill 

Suppose $R: G \rightarrow \GL_{n}(\R)$ is a representation of $G$ with respect to some basis $\mathscr{B}'$. 
Let 
$$\<x, y\>_{G} = \frac{1}{\#G} \sum_{g \in G} \<gx, gy\>,$$ 
where $\<x, y\>$ is the standard Hermitian, be a positive definite $G$-invariant Hermitian form. 
By the $\<, \>_G$ Gram Schmidt process, there is an orthonormal basis $\mathscr{B}$ for $\R^{n}$ with respect to the form $\<,\>_{G}$. 
Let $P$ be the basechange matrix from $\mathscr{B}'$ to $\mathscr{B}$. 
Take some $g \in G$ and let $A' = R_g$, that is, the matrix of $g$ with respect to the old basis $\mathscr{B}'$. 
The matrix $A$ of $g$ with respect to the new, orthonormal basis $\mathscr{B}$ is given by 
$$A = P^{-1} A' P.$$
Then one can see $A$ is unitary by checking $\<AX, AY\>_{G} = \<X, Y\>_{G}$ for all coordinate vectors $X, Y \in \R^{n}$ (with respect to $\mathscr{B}$). 
If $X = \bmat{x_1 \\ \vdots \\ x_n}$ and $Y = \bmat{y_1 \\ \vdots \\ y_n}$, 
\begin{align}
  \<AX, AY\>_{G} & = \left\<\sum_{j} x_j Ab_j, \sum_{k} y_k Ab_{k} \right\>_G \nonumber \\ 
                 & = \sum_{j, k} \bar{x}_j y_k \left\<Ab_j, Ab_k \right\>_G \nonumber \\ 
                 & = \sum_{j, k} \bar{x}_j y_k \left\<b_j, b_k \right\>_G \label{eq:5_G_invar} \\ 
                 & = \sum_{j = k} \bar{x}_j y_k = X^{*} Y = \<X, Y\>_G \label{eq:5_orthonorm}
\end{align}
where \eqref{eq:5_G_invar} follows because $Ab_j$ and $Ab_k$ are really the action of the group element $g$ on $b_j$ and $b_k$, respectively, and \eqref{eq:5_orthonorm} follows from the orthonormality of the basis $(b_1, \ldots, b_n)$.  
Hence $A$ is a unitary matrix; 
since the question restricts itself to $\GL_n(\R)$, rather than over the complex numbers, a unitary matrix is just an orthogonal matrix. 
Since this computation works for the basis $\mathscr{B}$ that was chosen independently from the group element $g \in G$, it follows that conjugating $R_g$ to represent it in the basis $\mathscr{B}$ gives an orthogonal matrix. 
So every finite subgroup $G$ of $\GL_{n}(\R)$ is conjugate to a subgroup of the orthogonal group $O_n$. 

%------------------------------------------------------------------------------%

\section{$G$-Invariant Hermitian Form for a Representation of $\mathbb{Z}/3\mathbb{Z}$}

Let $R \colon \mathbb{Z}/3\mathbb{Z} \to \mathrm{GL}_2(\mathbb{C})$ be the
representation with
\[
  R(\bar{1}) \;=\; \begin{pmatrix} -1 & -1 \\ 1 & 0 \end{pmatrix}.
\]
Construct a $G$-invariant Hermitian form on $\mathbb{C}^2$, then verify by
explicit matrix computation that your answer is correct.

\noindent\dotfill

Denote $G = Z_3$ and fix the standard basis on $\C^2$. 
Let $M \defeq R(\bar{1}) =  \bmat{-1 & -1 \\ 1 & 0}$. 
$M^2 = \bmat{0 & 1 \\ -1 & -1}$. 
Letting $\<, \>$ be the standard Hermitian product on $\C^2$, define 
\begin{align*}
  \<X, Y\>_R & = \frac{1}{\#G}\sum_{g \in G}\<g \cdot X, g \cdot Y\> \\ 
\end{align*}
where $X = \bmat{x_1 \\ x_2}, Y = \bmat{y_1 \\ y_2}$. 
The matrix of this $G$-invariant Hermitian form is given by $A = [a_{ij}]$, where $a_{ij} = \<e_i, e_j\>_R$.
For example, 
\begin{align*}
  a_{11} & = \<e_1, e_1\>_R \\ 
         & = \frac{1}{3} \left( \<I e_1, I e_1\> + \left\< \bmat{-1 & -1 \\ 1 & 0} e_1, \bmat{-1 & -1 \\ 1 & 0} e_1 \right\> + \left\< \bmat{0 & 1 \\ -1 & -1} e_1, \bmat{0 & 1 \\ -1 & -1} e_1 \right\> \right) \\ 
         & = \frac{1}{3}(1 + 2 + 1) \\ 
         & = \frac{4}{3}
\end{align*}
Calculating the other three values similarly, 
$$A = \frac{1}{3} \bmat{4 & 2 \\ 2 & 4}.$$ 
To verify that $A$ is indeed a $G$-invariant, one can calculate 
\begin{align*}
  M^{*}AM  & = A \\ 
  {M^2}^{*} A M^2 & = A, 
\end{align*}
hence for all $g \in G$, $\<g x, g y\>_R = \<x, y\>_R$. 

%------------------------------------------------------------------------------%

\section{A Non-Semisimple Representation in Characteristic $p$}

Show that the subgroup $G$ of $\mathrm{GL}_2(\mathbb{F}_p)$ generated by
$\begin{pmatrix} 1 & 1 \\ 0 & 1 \end{pmatrix}$
is cyclic of order $p$.  Then show that the vector space $\mathbb{F}_p^2$ does
admit a one-dimensional $G$-invariant subspace, but cannot be written as the
direct sum of two such subspaces.  (Note that this does not contradict
Maschke's theorem!)

\noindent\dotfill 

Let $R: G \rightarrow \GL_2(\F_p)$ be the representation in question. 
Notice that over $\Z$, 
$$\bmat{1 & 1 \\ 0 & 1}^n = \bmat{1 & n \\ 0 & 1}.$$
for any positive integer $n$. 
So clearly $M \defeq \bmat{1 & 1 \\ 0 & 1}$ is cyclic of order $p$ in $\GL_2(\F_p)$.
Since $R(G) \subset \GL_{2}(\F_p)$ is generated by this one element of order $p$, i.e., $R: 1_G \mapsto \bmat{1 & 1 \\ 0 & 1}$, one can conclude 
$$G \equiv \Z / p\Z.$$ 
To find the one-dimensional $G$-invariant subspace, let's determine the eigenvalue of $M$: 
one can quickly check that the characteristic polynomial of $M$ is $(t - 1)^2$. 
Hence $\lambda = 1$ is the only eigenvalue of $M$, and its only corresponding eigenvector (up to scaling) is $e_1$:
$$\bmat{1 & 1 \\ 0 & 1}\bmat{1 \\ 0} = \bmat{1 \\ 0}.$$
So $W \defeq \text{span}(e_1) \subset \F_{p}^{2}$ is a one-dimensional $G$-invariant subspace. 

It is clear that $W^{\perp} = \text{span}(e_2)$, where perp is taken with respect to the standard Hermitian.  
But clearly $e_2$ is not invariant under $G$: 

However, no one-dimensional subspace of $\F_{p}^{2}$ linearly independent from $W$ is $G$-invariant. 
If there was, then that subspace would be generated by vector $v \in \F_{p}^{2}$ linearly independent from $e_1$. 
By $G$-invariance, $Mv = \lambda v$ for some $\lambda$, i.e., $v$ must be a eigenvector of $M$. 
But $M$ only has eigenvalue 1, whose eigenvector, up to scaling, is $e_1$, forcing $v \in W$, a contradiction. 
Hence $\F_{p}^{2}$ cannot be written as the direct product of two one-dimensional $G$-invariant subspaces. 

%------------------------------------------------------------------------------%

\section{Nonabelian Group of Order 55}

Suppose that $G$ is a nonabelian group of order $55$.  Compute the class
equation and the dimensions of the irreducible characters of $G$.

\emph{Hint:} you should be able to do this just from the known constraints on
conjugacy classes and characters.  For an extra challenge, you could try to
prove that this group exists and is unique, and then use its structure to read
off the requested information.

\noindent\dotfill 

Let $n_5$ and $n_{11}$ denote the number of Sylow-5 and Sylow-11 subgroups of $G$, respectively. 
By Sylow's third theorem, 
\begin{align*}
  n_{11} = 1 \pmod{11}, \hspace{1.5em} & n_{11} \mid 5 \\ 
  n_{5} = 1 \pmod{5}, \hspace{1.5em} & n_{5} \mid 11 
\end{align*}
$n_{11} = 1$ follows immediately, whereas there are two possible values for $n_5$: 1 or 11. 
If $n_5 = 1$, then by Sylow's Second Theorem, the unique Sylow-5 and -11 subgroups $P_5$ and $P_{11}$ are both normal in $G$. 
Notice that $P_{5}$ and $P_{11}$ have prime orders, so necessarily they are cyclic groups of order 5 and 11, respectively. 
By the fundamental theorem for product groups or whatever it's called, $G \equiv Z_{5} \times Z_{11}$, which contradicts the assumption that $G$ is nonabelian. 
Hence $n_{5} = 11$. 

At this point, I know that $G$ is the semidirect product of something, but I really don't know that theory too well; 
let's do something else. 
Notice that given any two $P_{5} \neq P_{5}' \in \Syl{5}{G}$, $P_5 \cap P_{5}' = \{1\}$ (if not, the shared nontrivial element generates both $P_{5}$ and $P_{5}'$ because they are both isomorphic to $Z_5$, and hence they are canonically equal subgroups). 
So $P_{11}$ and each of the $P_5$'s are pairwise disjoint besides the identity element. 
Counting the distinct elements of $P_{11}$ and the eleven $P_5$'s, 
\begin{equation} \label{eq:8_syl_sum}
  1 + 10 + 11(4) = 55.
\end{equation}
where the 1 in the left hand side corresponds to the identity element in $G$. 
In other words, each nontrivial element of $G$ belongs to exactly one Sylow subgroup of $G$. 

Now, let's try to decompose the numbers in \eqref{eq:8_syl_sum} to obtain $G$'s class equation. 
First, each number in the class equation must divide $\#G$; 
in this case, this means each number is either 1, 5, or 11. 
Recall that any normal subgroup is the union of conjugacy classes. 
The unique $P_{11}$ is normal by Sylow's second theorem. 
Since the 10 in the left hand side of \eqref{eq:8_syl_sum} corresponds to $P_{11}$, the class equation will partially simplify to 
\begin{equation} \label{eq:8_simplify}
  1 + 10 + c_1 + \ldots + c_k = 55
\end{equation}
where 10 is the sum of the sizes of the conjugacy classes of nontrivial elements in $P_{11}$, 
and where $c_1, \ldots, c_k$ are the sizes of the conjugacy classes of the elements not in $P_{11}$; 
that is, sizes of conjugacy classes of the nontrivial elements in any of the $P_5$'s. 
Notice that none of the $c_i = 1$. 
If some $c_i = 1$, then some $P_5$ contains a nonidentity element in $Z(G)$, the center of $G$. 
That element generates the $P_5$, so the entire $P_5$ is contained in $Z(G)$. 
Namely, then the $P_5$ is normal. 
This contradicts Sylow's second theorem and the fact that $n_5 = 11 \neq 1$  
(in general, $P_p$ is normal in $G$ if and only if $n_p = 1$). 
So the conclusion is that the $c_i$'s are either 5 or 11. 
In fact, considering \eqref{eq:8_simplify} gives $c_1 + \ldots + c_k = 44$, basic arithmetic shows that $k = 4$ and each $c_{i} = 11$. 

The only thing left to determine is the conjugacy classes that make up the unique $P_{11}$. 
Since $10 \nmid 55$, it is clear $P_{11} \setminus \{1_G\}$ is itself not a conjugacy class. 
I show that $P_{11}$ does not contain any nontrivial element also belonging to $Z(G)$; 
that is, there is no conjugacy class of size 1, besides the conjugacy class of the identity, contained in $P_{11}$.  
It was demonstrated in the previous paragraph that none of the 44 nonidentity elements in a $P_5$ is in $Z(G)$. 
Because $Z(G)$ is a (normal) subgroup of $G$, it follows by considering the order of $G$ that $Z(G) = \{1_G\}$ or $Z(G) = P_{11}$. 
Suppose $Z(G) = P_{11}$. 
Then $\#(G / Z(G)) = 5$, i.e., $G / Z(G)$ is cyclic. 
Recall that if $G / Z(G)$ is cyclic, then $G$ is abelian. 
This contradicts the assumption that $G$ is nonabelian. 
Hence $Z(G) = \{1\}$ and $P_{11}$ contains no nontrivial element that is in the center. 
This forces $P_{11} \setminus \{1\}$ to be the union of two conjugacy classes of size five. 
In all,  
$$1 + 5 + 5 + 11 + 11 + 11 + 11 = 55$$ 
is the class equation for $G$. 

Determining the irreducible characters of $G$ follows quickly. 
Since $G$ has seven conjugacy classes, there are seven isomorphism classes of the irreducible representations of $G$. 
Letting $d_1, \ldots, d_7$ be their dimensions, there are two restrictions: 
\begin{enumerate}
  \item $d_{1}^{2} + \ldots + d_{7}^{2} = 55$, 
  \item Each $d_i \mid 55$. 
\end{enumerate}
It is clear from combining the two conditions that each $d_i \in \{1, 5\}$. 
With this restraint, the only tuple that satisfies condition 1 is 
$$(d_1, d_2, d_3, d_4, d_5, d_6, d_7) = (1, 1, 1, 1, 1, 5, 5),$$
which finishes this problem.


\end{document}
